% !TeX program = xelatex
% Compile twice: xelatex triangular_commutator_combined.tex
\documentclass[UTF8,11pt,a4paper,fontset=fandol]{ctexart}
\usepackage[margin=26mm]{geometry}
\usepackage{amsmath,amssymb,amsthm,mathtools}
\usepackage{xcolor}
\usepackage[colorlinks=true,linkcolor=blue!60!black,citecolor=blue!60!black,urlcolor=blue!60!black]{hyperref}
\usepackage{enumitem}
\setlist[enumerate]{itemsep=3pt,topsep=4pt}
\setlength{\emergencystretch}{2em}
\newtheorem{theorem}{定理}[section]
\newtheorem{lemma}[theorem]{引理}
\newtheorem{corollary}[theorem]{推论}
\theoremstyle{remark}
\newtheorem{remark}[theorem]{说明}
\newtheorem{example}[theorem]{例}
\renewcommand{\proofname}{证明}
\newcommand{\BH}{B(H)}
\newcommand{\Sp}{S_r(H)}
\newcommand{\Tr}{\operatorname{Tr}}
\newcommand{\dist}{\operatorname{dist}}
\newcommand{\PP}{\mathcal P}
\newcommand{\QQ}{\mathcal Q}
\newcommand{\DD}{\mathcal D}
\newcommand{\TriH}{\mathcal H}
\newcommand{\Alg}{\mathcal A}
\newcommand{\CC}{\mathcal C}
\newcommand{\norm}[1]{\left\lVert #1\right\rVert}
\DeclareMathOperator{\rank}{rank}
\newcommand{\HS}{\mathfrak H}
\numberwithin{equation}{section}
\title{三角投影交换子的有界性与 Schatten 类刻画\\
\large 关于 Power (1985) 定理 2.1 的两个补充结果}
\author{}
\date{}
\begin{document}
\maketitle
\vspace{-2em}

本文在同一个离散有限秩巢的框架下，统一讨论左乘交换子
$\CC_X=[L_X,\PP]$ 的两个性质：
\begin{enumerate}[label=\textup{(\arabic*)}]
 \item $\CC_X$ 在非交换 $L_r(B(H),\Tr)=S_r(H)$ 上有界，$1<r<\infty$；
 \item 将 $\CC_X$ 看成 Hilbert 空间 $S_2(H)$ 上的算子时，
       $\CC_X\in S_p(S_2(H))$，$0<p<\infty$。
\end{enumerate}
为区分这两个空间的指标，第一类问题使用 $r$，第二类问题使用 $p$。
第一类问题中，原文已有的假设 $X\in B(H)$ 足以保证有界性；
若允许形式块矩阵符号，则充要条件为两个方向的角块一致有界。
第二类问题的充要条件则是角块的 Schatten 范数满足加权可和性。
文中保留完整证明、自伴情形、Hilbert--Schmidt 特例和相关反例，
并在附录证明三角截断的 $S_r$ 有界性。

\section{统一设置与有界符号的直接结论}

设 $H$ 为可分复 Hilbert 空间，
\[
 0=P_0<P_1<P_2<\cdots,\qquad \rank P_n<\infty,
 \qquad P_n\longrightarrow I\quad\text{强算子收敛},
\]
并令 $E_n=P_n-P_{n-1}$，记 $d_n=\rank E_n$、$\HS=S_2(H)$。原文的巢代数为
\[
 \Alg=\{T\in\BH:(I-P_n)TP_n=0\ \text{对所有 }n\},
\]
即相对于 $H=\bigoplus_{n\geq1}E_nH$ 的块上三角代数。
在有限支撑块矩阵上定义
\[
 \PP(A)=\sum_{i\leq j}E_iAE_j,\qquad
 \QQ=I-\PP,\qquad
 \DD(A)=\sum_jE_jAE_j.
\]
当 $r=2$ 时，$\PP$ 就是原文投影到 $\Alg\cap S_2$ 的正交投影。
对每个 $1<r<\infty$，三角截断定理给出其唯一有界延拓，并且
\begin{equation}\label{eq:tp}
 t_r:=\norm{\PP}_{S_r\to S_r}<\infty,\qquad
 q_r:=\norm{\QQ}_{S_r\to S_r}<\infty.
\end{equation}
这些常数可以只依赖于 $r$，与块的个数和维数无关。
为使证明完整，附录给出这一事实的一个证明。

\medskip
原文 $X\PP-\PP X$ 中的 $X$ 表示\emph{左乘算子}
$L_X(A)=XA$。因此本文研究的是
\begin{equation}\label{eq:defC}
 \CC_X=[L_X,\PP],\qquad
 \CC_X(A)=X\PP(A)-\PP(XA),
\end{equation}
它与内部导子 $A\mapsto XA-AX$ 是不同的映射。

\begin{theorem}[保留原文的有界符号假设]\label{thm:bounded}
设 $X\in\BH$。则对每个 $1<r<\infty$，
$\CC_X$ 是 $S_r(H)$ 上的有界线性算子，且
\begin{equation}\label{eq:easybound}
 \norm{\CC_X}_{S_r\to S_r}\leq 2t_r\norm{X}_{\BH}.
\end{equation}
特别地，不要求 $X$ 紧、不要求 $X\in S_r$，也不要求
$X\in(\Alg+K(H))\cap(\Alg+K(H))^*$。
\end{theorem}

\begin{proof}
Schatten 理想不等式给出
\[
 \norm{XA}_r\leq\norm{X}_{\BH}\norm{A}_r,
 \qquad A\in S_r(H).
\]
结合 \eqref{eq:tp}，有
\begin{align*}
 \norm{X\PP(A)-\PP(XA)}_r
 &\leq\norm{X}_{\BH}\norm{\PP(A)}_r
       +t_r\norm{XA}_r\\
 &\leq 2t_r\norm{X}_{\BH}\norm{A}_r.
\end{align*}
这证明了结论。原文定理 2.1(ii) 中的代数条件刻画的是
$S_2$ 上的\emph{紧性}，而上述结论只涉及有界性。
\end{proof}

\section{若不预设 \texorpdfstring{$X$}{X} 有界：精确的符号条件}

本节的强化结论适用于原文的同一离散巢，包括有限维块的重数情形。
允许 $X=(X_{ij})_{i,j\geq1}$ 仅为形式块矩阵，
其中 $X_{ij}:E_jH\to E_iH$；此时不预设 $X$ 在 $H$ 上定义有界算子。
令 $X^{(N)}=(X_{ij})_{1\leq i,j\leq N}$，视为 $P_NH$ 上的算子，并定义
\begin{align}
 \beta_-(X)&=\sup_{N>n\geq1}
     \norm{(P_N-P_n)X^{(N)}P_n},\label{eq:beta-}\\
 \beta_+(X)&=\sup_{N>n\geq1}
     \norm{P_nX^{(N)}(P_N-P_n)},\label{eq:beta+}\\
 \beta(X)&=\max\{\beta_-(X),\beta_+(X)\}.\label{eq:beta}
\end{align}
这里所有范数均为 Hilbert 空间算子范数。
这些有限截断的定义避免了无界矩阵乘法的定义域问题。
若两个方向的角块已定义为有界算子，则等价地
\begin{equation}\label{eq:cut}
 \beta(X)=\sup_n\max\{
 \norm{(I-P_n)XP_n},\norm{P_nX(I-P_n)}\}
 =\sup_n\norm{[X,P_n]}.
\end{equation}
最后一个等号来自相对于 $P_nH\oplus(I-P_n)H$ 的表示
\[
 [X,P_n]=\begin{pmatrix}
 0&-P_nX(I-P_n)\\
 (I-P_n)XP_n&0
 \end{pmatrix}.
\]

对有限支撑 $A=(A_{kj})$，交换子的形式定义为
\begin{equation}\label{eq:entries}
 (\CC_X A)_{ij}
 =\sum_k X_{ik}A_{kj}
 \bigl(\mathbf1_{\{k\leq j\}}-\mathbf1_{\{i\leq j\}}\bigr).
\end{equation}
每个块的求和有限，但所得整个矩阵未必属于 $S_r$。
所谓其在 $S_r$ 上有界，是指存在 $S_r\to S_r$ 的有界映射，
在有限支撑矩阵上逐块满足 \eqref{eq:entries}。

\begin{theorem}[一致角块条件]\label{thm:sharp}
对上述形式块矩阵 $X$，下列条件等价：
\begin{enumerate}[label=\textup{(\roman*)}]
 \item 对某个 $1<r<\infty$，$\CC_X$ 在 $S_r(H)$ 上有界；
 \item $\beta(X)<\infty$；
 \item 对所有 $1<r<\infty$，$\CC_X$ 在 $S_r(H)$ 上有界。
\end{enumerate}
而且对每个 $1<r<\infty$，有
\begin{equation}\label{eq:sharpbounds}
 \boxed{\quad
 \beta(X)\leq\norm{\CC_X}_{S_r\to S_r}
 \leq t_pq_r\bigl(\beta_-(X)+\beta_+(X)\bigr)
 \leq2t_pq_r\beta(X).
 \quad}
\end{equation}
在 $r=2$ 时，还有精确等式
\begin{equation}\label{eq:exact2}
 \norm{\CC_X}_{S_2\to S_2}=\beta(X).
\end{equation}
\end{theorem}

先列出证明所需的有界完成引理。它使用原文在定理 2.1
之前已证明的距离公式
\begin{equation}\label{eq:nehari}
 \norm{\QQ L_Y\PP}_{S_2\to S_2}=\dist(Y,\Alg),
 \qquad Y\in\BH;
\end{equation}
见 \cite[p.~277]{Power}，其依据为原文定理 1.6 及其后的说明。
这里的距离取 $B(H)$ 算子范数。

\begin{lemma}[严格三角部分的有界完成]\label{lem:completion}
若 $\beta_-(X)<\infty$，则存在 $B\in\BH$ 满足
\[
 E_iBE_j=X_{ij}\quad(i>j),\qquad
 \norm B\leq\beta_-(X).
\]
若 $\beta_+(X)<\infty$，则存在 $C\in\BH$ 满足
\[
 E_iCE_j=X_{ij}\quad(i<j),\qquad
 \norm C\leq\beta_+(X).
\]
注意这里是为给定的严格三角部分\emph{补齐其他块}，
并未断言将其余块置零后所得的三角矩阵本身有界。
\end{lemma}

\begin{proof}
先说明对任意 $Y\in\BH$，
\begin{equation}\label{eq:Hcorner}
 \norm{\QQ L_Y\PP}_{S_2\to S_2}
 =\sup_j\norm{(I-P_j)YP_j}.
\end{equation}
事实上，逐块列计算得到
\[
 (\QQ L_Y\PP)(A)E_j=(I-P_j)YP_jAE_j.
\]
由 $\norm A_2^2=\sum_j\norm{AE_j}_2^2$ 得到 ``$\leq$''。
固定 $j$，取 $A=\xi\otimes\eta^*$，其中 $\xi\in P_jH$、
$\eta\in E_jH$ 均为单位向量，再对 $\xi$ 取上确界，得到反向不等式。
这里 $\xi\otimes\eta^*$ 表示 $h\mapsto\langle h,\eta\rangle\xi$。

将 $Y_N=X^{(N)}$ 在 $(I-P_N)H$ 上补零。
由 \eqref{eq:Hcorner}、\eqref{eq:nehari} 以及 \eqref{eq:beta-}，有
\[
 \dist(Y_N,\Alg)\leq\beta_-(X).
\]
故可取 $T_N\in\Alg$ 使
\[
 \norm{Y_N-T_N}\leq\beta_-(X)+N^{-1}.
\]
令 $B_N=P_N(Y_N-T_N)P_N$。则 $\norm{B_N}\leq\beta_-(X)+N^{-1}$，
且当 $N\geq i>j$ 时，$E_iB_NE_j=X_{ij}$。
由于 $H$ 可分，有界序列 $(B_N)$ 有弱算子收敛子列。
设其极限为 $B$，便有 $\norm B\leq\beta_-(X)$，
且每个严格下三角块均保持为 $X_{ij}$。

对形式伴随矩阵 $X^*=(X_{ji}^*)$ 应用刚证明的结论，
并使用 $\beta_-(X^*)=\beta_+(X)$，再取伴随，即得到 $C$。
\end{proof}

\begin{proof}[定理 \ref{thm:sharp} 的证明]
\textbf{必要性。}
设 $\CC_X$ 在某个 $S_r$ 上有界，记其范数为 $M$。
固定 $N>n$，取单位向量 $\eta\in E_nH$ 及任意 $\xi\in P_NH$，
令 $A=\xi\otimes\eta^*$。由 \eqref{eq:entries} 可得
\[
 P_N\CC_X(A)P_N
 =\bigl([X^{(N)},P_n]\xi\bigr)\otimes\eta^*.
\]
秩一算子的每个 Schatten 范数均等于两向量范数之积，
而压缩在 $S_r$ 上是压缩映射，故
\[
 \norm{[X^{(N)},P_n]\xi}
 \leq\norm{\CC_X(A)}_r\leq M\norm\xi.
\]
于是 $\norm{[X^{(N)},P_n]}\leq M$。
该范数正是两个相反角块范数的最大值。
对 $N,n$ 取上确界便得 $\beta(X)\leq M$。

\medskip
\textbf{充分性。}
设 $\beta(X)<\infty$。由引理 \ref{lem:completion} 取有界完成 $B,C$。
定义 $S_r$ 上的有界映射
\begin{equation}\label{eq:factorization}
 T=\QQ L_B\PP-\PP L_C\QQ.
\end{equation}
对有限支撑矩阵，直接计算得到
\begin{align*}
 (\QQ L_B\PP A)_{ij}
 &=\begin{cases}
     \displaystyle\sum_{k\leq j}X_{ik}A_{kj},&i>j,\\
     0,&i\leq j,
   \end{cases}\\
 (\PP L_C\QQ A)_{ij}
 &=\begin{cases}
     \displaystyle\sum_{k>j}X_{ik}A_{kj},&i\leq j,\\
     0,&i>j.
   \end{cases}
\end{align*}
其中第一式仅用到 $i>j\geq k$ 时 $B_{ik}=X_{ik}$，
第二式仅用到 $i\leq j<k$ 时 $C_{ik}=X_{ik}$。
因此 $T$ 逐块等于 \eqref{eq:entries}，从而就是所需交换子的延拓。
有限支撑矩阵在 $S_r$ 中稠密，故延拓唯一。
此外，
\[
 \norm T_{S_r\to S_r}
 \leq q_pt_r(\norm B+\norm C)
 \leq t_pq_r\bigl(\beta_-(X)+\beta_+(X)\bigr).
\]
这对每个 $1<r<\infty$ 均成立，证明了等价性和 \eqref{eq:sharpbounds}。

\medskip
\textbf{$S_2$ 上的等式。}
一致有界的角块定义了有界算子 $K_j=[X,P_j]$，且
$\norm{K_j}\leq\beta(X)$。
逐块列有
\[
 \CC_X(A)E_j=K_jAE_j.
\]
所以
\[
 \norm{\CC_X(A)}_2^2
 =\sum_j\norm{K_jAE_j}_2^2
 \leq\beta(X)^2\sum_j\norm{AE_j}_2^2
 =\beta(X)^2\norm A_2^2.
\]
结合已证的下界即得 \eqref{eq:exact2}。
\end{proof}

\begin{remark}[条件的含义及其与原文定理 2.3 的关系]
\leavevmode
\begin{enumerate}[label=\textup{(\arabic*)}]
 \item $\beta$ 是一个忽略块对角部分的半范数。
 若 $D=(D_{ij})$ 为形式块对角矩阵，即 $D_{ij}=0$ 对 $i\ne j$，则
 \[
  \CC_{X+D}=\CC_X,\qquad\beta(X+D)=\beta(X).
 \]
 因而允许形式符号时，不能把 ``$X\in B(H)$'' 当作必要条件。
 \item 在原文定理 2.3 所考虑的重数为一情形，将对角线归零以后，
 $\beta(X)<\infty$ 等价于
 \[
 X-\DD X\in B(H)+\widetilde{B(H)},
 \]
 其中 $\widetilde Y=-i\TriH Y$ 是原文的三角共轭，
 $(\TriH Y)_{ij}=\operatorname{sgn}(j-i)Y_{ij}$。
 这正是原文所讨论的 BMO 型符号条件；不能在未指定定义时
 将它直接替换为其他非交换 BMO 空间。
 \item 本文的充要条件针对原文的离散有限秩巢和左乘交换子。
 它没有把问题改成一般半有限 von Neumann 代数上的任意投影问题。
\end{enumerate}
\end{remark}

\section{作为 Hilbert--Schmidt 空间上算子的 Schatten 类刻画}

本节先取 $X\in B(H)$，并沿用第一节的 $E_n,d_n$ 和 $\HS=S_2(H)$。
$S_p(\HS)$ 指 Hilbert 空间 $\HS$ 上算子的 Schatten 类；
其定义使用 $\CC_X$ 本身的奇异值，而非 $\CC_X:S_r(H)\to S_r(H)$ 的算子范数。
第 \ref{sec:compact} 节还将说明，若形式符号已产生 $S_2(H)$ 上的有界交换子，同一公式仍然成立。

\subsection{充要条件与精确范数公式}

\begin{theorem}\label{sch:thm:main}
对任意 $0<p<\infty$（特别是 $1<p<\infty$），下列条件等价：
\[
 \boxed{\quad
 \CC_X\in S_p\bigl(S_2(H)\bigr)
 \quad\Longleftrightarrow\quad
 \sum_{n=1}^{\infty}d_n\norm{[X,P_n]}_{S_p(H)}^p<\infty.
 \quad}
\]
而且有精确等式
\begin{equation}\label{sch:eq:main}
 \norm{\CC_X}_{S_p(S_2(H))}^p
 =\sum_{n=1}^{\infty}d_n\norm{[X,P_n]}_{S_p(H)}^p.
\end{equation}
等价地，该条件可以写为
\begin{equation}\label{sch:eq:corners}
 \boxed{\quad
 \sum_{n=1}^{\infty}d_n\left(
 \norm{(I-P_n)XP_n}_{S_p(H)}^p
 +\norm{P_nX(I-P_n)}_{S_p(H)}^p\right)<\infty.
 \quad}
\end{equation}
\eqref{sch:eq:corners} 中的和也精确等于 \eqref{sch:eq:main} 的左侧。
对于 $0<p<1$，这里使用 Schatten 拟范数。
\end{theorem}

\begin{proof}
\textbf{第一步：将交换子分解为角块交换子的直和。}
对每个 $n$，取 $E_nH$ 的一组标准正交基
$\{\eta_{n,r}:1\leq r\leq d_n\}$。
所有这些向量合起来构成 $H$ 的标准正交基。
由 Hilbert--Schmidt 范数的列表示，映射
\[
 U:\HS\longrightarrow\bigoplus_{n\geq1}\bigoplus_{r=1}^{d_n}H,
 \qquad U(A)=(A\eta_{n,r})_{n,r}
\]
是酉算子，因为
\[
 \norm A_2^2=\sum_{n\geq1}\sum_{r=1}^{d_n}\norm{A\eta_{n,r}}_H^2,
\]
且任意平方可和的向量列均唯一确定一个 Hilbert--Schmidt 算子。

固定 $n,r$。由于 $\eta_{n,r}\in E_nH$，有
\[
 \PP(A)\eta_{n,r}=P_nA\eta_{n,r}.
\]
因此
\begin{align*}
 \CC_X(A)\eta_{n,r}
 &=XP_nA\eta_{n,r}-P_nXA\eta_{n,r}\\
 &=[X,P_n]A\eta_{n,r}.
\end{align*}
于是得到了真正的酉等价
\begin{equation}\label{sch:eq:unitary}
 \boxed{\quad
 U\CC_XU^{-1}
 =\bigoplus_{n\geq1}\bigoplus_{r=1}^{d_n}[X,P_n].
 \quad}
\end{equation}
换言之，$[X,P_n]$ 的每个非零奇异值，在这一分量中恰好重复 $d_n$ 次。

\medskip
\textbf{第二步：对奇异值求 $p$ 次幂和。}
每个 $K_n=[X,P_n]$ 都是有限秩算子，因为 $P_n$ 有限秩且 $X$ 有界。
直和算子的奇异值由各分量的非零奇异值汇集而成，只要这些奇异值趋于零。
若 $\sum_nd_n\norm{K_n}_p^p<\infty$，则
\[
 \norm{K_n}^p\leq\norm{K_n}_p^p\longrightarrow0.
\]
因此 \eqref{sch:eq:unitary} 的直和是紧算子，其奇异值 $p$ 次幂之和正是
$\sum_nd_n\norm{K_n}_p^p$，证明了充分性及公式 \eqref{sch:eq:main}。
反过来，若 $\CC_X\in S_p(\HS)$，
则其有限分量直和的奇异值 $p$ 次幂和一致有界；
对分量数取极限，即得同一公式和该级数的收敛性。
这一论证对所有 $p>0$ 都成立。

\medskip
\textbf{第三步：改写为两个方向的角块条件。}
设
\[
 B_n=(I-P_n)XP_n,
 \qquad A_n=P_nX(I-P_n).
\]
相对于 $H=P_nH\oplus(I-P_n)H$，有
\[
 K_n=[X,P_n]=\begin{pmatrix}0&-A_n\\B_n&0\end{pmatrix},
 \qquad
 K_n^*K_n=\begin{pmatrix}B_n^*B_n&0\\0&A_n^*A_n\end{pmatrix}.
\]
从而对任意 $p>0$，
\[
 \norm{K_n}_p^p=\norm{B_n}_p^p+\norm{A_n}_p^p.
\]
代入 \eqref{sch:eq:main}，即得 \eqref{sch:eq:corners}。
\end{proof}

\begin{remark}[为何会出现重数 $d_n$]
在 $S_2(H)E_n$ 上，左乘 $[X,P_n]$ 对 $E_nH$ 的每一个基向量所对应的列，
都独立作用一次。因此一个块列包含 $d_n$ 个相同的 Hilbert 空间算子分量。
重数为一时 $d_n=1$；一般的有限维块情形不能省略这个因子。
\end{remark}

\begin{corollary}[自伴符号]
若 $X=X^*$，则两个角块互为伴随，因而
\[
 \norm{\CC_X}_{S_p(S_2(H))}^p
 =2\sum_{n\geq1}d_n\norm{(I-P_n)XP_n}_{S_p(H)}^p.
\]
所以自伴情形只需检验一个方向的角块的加权 $p$ 次幂可和性。
\end{corollary}

\subsection{重数为一的矩阵形式及 Hilbert--Schmidt 特例}

在标准巢 $H=\ell_2(\mathbb N)$、
$P_n$ 投影到 $\operatorname{span}\{e_1,\ldots,e_n\}$ 的情形，
令 $X=(x_{ij})$，则 $d_n=1$，上述条件成为
\begin{equation}\label{sch:eq:matrix}
 \sum_{n=1}^{\infty}\left(
  \norm{(x_{ij})_{i>n,\,j\leq n}}_{S_p}^p
  +\norm{(x_{ij})_{i\leq n,\,j>n}}_{S_p}^p
 \right)<\infty.
\end{equation}
这里的矩阵表示相应两个 Hilbert 空间之间的角块算子，
范数为其奇异值的 Schatten 范数。

当 $p=2$ 时可以把这一条件完全展开：
\begin{align}
 \norm{\CC_X}_{S_2(S_2(H))}^2
 &=\sum_{n\geq1}\left(
   \sum_{i>n,\,j\leq n}|x_{ij}|^2
   +\sum_{i\leq n,\,j>n}|x_{ij}|^2\right)\notag\\
 &=\boxed{\ \sum_{i\ne j}|i-j|\,|x_{ij}|^2\ }.\label{sch:eq:hs}
\end{align}
事实上，当 $i>j$ 时，$x_{ij}$ 恰好出现在 $n=j,j+1,\ldots,i-1$
这 $i-j$ 个下角块中；当 $i<j$ 时，计数同理。
所以
\[
 \CC_X\in S_2(S_2(H))
 \quad\Longleftrightarrow\quad
 \sum_{i\ne j}|i-j|\,|x_{ij}|^2<\infty.
\]
对于 $p\ne2$，应使用 \eqref{sch:eq:matrix} 中的角块奇异值范数；
上述计算不提供将其改成逐项 $|x_{ij}|^p$ 加权和的依据。

更一般地，在有限维块情形，令 $r_n=\rank P_n$、$r_0=0$，则
\begin{equation}\label{sch:eq:blockhs}
 \norm{\CC_X}_{S_2(S_2(H))}^2
 =\sum_{i\ne j}|r_{i-1}-r_{j-1}|\,
   \norm{E_iXE_j}_{S_2(H)}^2.
\end{equation}
这是因为每个 $E_iXE_j$ 的权重为
$\sum_{n=\min(i,j)}^{\max(i,j)-1}d_n$。

\section{紧性、反例及形式符号的说明}\label{sec:compact}

由 \eqref{sch:eq:unitary}，还立即得到
\[
 \CC_X\text{ 在 }S_2(H)\text{ 上紧}
 \quad\Longleftrightarrow\quad
 \norm{[X,P_n]}_{B(H)}\longrightarrow0.
\]
这是原文定理 2.1(ii) 的紧性条件的角块表述。
本题要求的 Schatten 条件则进一步要求角块的\emph{全部奇异值}满足
\eqref{sch:eq:main} 的加权可和性。
本文的直和分解与 \cite[p.~279, Remark (2)]{Power} 对 Hankel 分量的
分解相呼应，但 \eqref{sch:eq:unitary} 是直接针对整个交换子得到的公式。

\begin{example}[即使 $X$ 为秩一，也不保证交换子属于给定 Schatten 类]
固定 $p>0$，在标准重数为一巢上令
\[
 u_i=i^{-1/2-1/p},\qquad
 X=u\otimes e_1^*.
\]
由于 $u\in\ell_2$，$X$ 是有界秩一算子，属于每个 $S_q(H)$，$q>0$。
然而对 $n\geq1$，有
\[
 P_nX(I-P_n)=0,\qquad
 (I-P_n)XP_n=((I-P_n)u)\otimes e_1^*.
\]
所以
\[
 \norm{[X,P_n]}_{S_p(H)}^p
 =\left(\sum_{i>n}i^{-1-2/p}\right)^{p/2}
 \asymp\frac1n.
\]
由调和级数发散及定理 \ref{sch:thm:main}，
$\CC_X\notin S_p(S_2(H))$。
同时角块算子范数趋于零，故这个交换子仍然是紧算子。
这说明条件不能简化为 $X\in S_p(H)$，甚至 $X$ 为有限秩也不够。
\end{example}

\begin{remark}[对角部分不影响结论]
若 $D\in B(H)$ 为相对于 $(E_n)$ 的块对角算子，则
$[D,P_n]=0$，因此 $\CC_{X+D}=\CC_X$。
本题的条件和范数均只取决于 $X$ 的非对角块。
\end{remark}

\begin{remark}[只假定交换子有界时仍然成立]
若采用定理 \ref{thm:sharp} 中允许形式块矩阵符号的设置，只假定
$\CC_X$ 已延拓为 $S_2(H)$ 上的有界算子，定理仍然成立。
事实上，以 $\xi\otimes\eta_{n,r}^*$ 作输入，逐列公式表明形式矩阵
$[X,P_n]$ 唯一确定一个范数不超过 $\norm{\CC_X}$ 的有界算子。
它相对于 $P_nH\oplus(I-P_n)H$ 只有两个非对角角块，故仍然为有限秩。
于是 \eqref{sch:eq:unitary} 及其后的奇异值计算完全适用。
因此公式要求的是角块算子的有界性及可和性，而不依赖于对角符号是否有界。
\end{remark}

\subsection{有界性与紧性的例子}

\begin{example}[有界交换子可以不紧]
在 $H=\ell_2(\mathbb N)$ 的标准重数为一巢上，令
$Se_n=e_{n+1}$ 为单边移位。由于 $S\in B(H)$，定理 \ref{thm:bounded}
给出 $\CC_S$ 在所有 $1<r<\infty$ 的 $S_r$ 上有界。
但对矩阵单位 $e_{nn}$，有
\[
 \CC_S(e_{nn})=e_{n+1,n}.
\]
这些输入和输出的 $S_r$ 范数均为 $1$，且不同输出之间的距离为
$2^{1/r}$，故 $\CC_S$ 不紧。
\end{example}

\begin{example}[无界对角符号也可以产生有界交换子]
仍在标准巢上，取稠密定义的对角算子 $De_n=ne_n$。
其有限支撑矩阵交换子恒为零，故唯一延拓是 $S_r$ 上的零映射。
这说明精确的符号条件必须忽略对角部分。
此例按有限支撑定义交换子，不主张两个含无界乘子的项
分别对每个 $A\in S_r$ 都有定义。
\end{example}

\appendix
\section{三角截断在 \texorpdfstring{$S_r$}{Sr} 上有界的证明}

这里给出 \eqref{eq:tp} 的证明，所用基础工具只有 Schatten H\"older
不等式、对偶性和复插值。先在任意有限个有限维块上工作，
所有估计均与块的数目及维数无关。
设
\[
 (\TriH A)_{ij}=\operatorname{sgn}(j-i)A_{ij},\qquad
 \operatorname{sgn}(0)=0.
\]
块对角映射 $\DD$ 在每个 $S_r$（$1\leq r\leq\infty$）上均为压缩映射：
它是块对角标量酉算子的共轭作用的平均。
此外，$\norm{\TriH}_{S_2\to S_2}\leq1$。

\medskip
对任意块矩阵 $A,B$，逐项核对符号可得 Cotlar 型恒等式
\begin{equation}\label{eq:cotlar}
 (\TriH A)(\TriH B)
 =\TriH\bigl((\TriH A)B+A(\TriH B)\bigr)-AB+(\DD A)(\DD B).
\end{equation}
其依据是对任意 $i,k,j$ 有
\begin{align*}
 &\operatorname{sgn}(k-i)\operatorname{sgn}(j-k)\\
 &\quad=\operatorname{sgn}(j-i)
   \bigl(\operatorname{sgn}(k-i)+\operatorname{sgn}(j-k)\bigr)
    -1+\mathbf1_{\{i=k=j\}}.
\end{align*}
由于 $\TriH(A^*)=-(\TriH A)^*$，在 \eqref{eq:cotlar} 中以 $A^*,A$
替代 $A,B$ 得到
\begin{equation}\label{eq:square}
 (\TriH A)^*(\TriH A)
 =A^*A-(\DD A)^*(\DD A)
   +\TriH\bigl((\TriH A)^*A-A^*\TriH A\bigr).
\end{equation}
设已经知道 $\norm{\TriH Z}_r\leq h_r\norm Z_r$，其中 $h_r$ 与维数无关。
对 \eqref{eq:square} 使用 $S_r$ 范数及 H\"older 不等式，可得
\begin{align*}
 \norm{\TriH A}_{2r}^{2}
 &\leq 2\norm A_{2r}^{2}
       +2h_r\norm{\TriH A}_{2r}\norm A_{2r}.
\end{align*}
解这一关于 $\norm{\TriH A}_{2r}/\norm A_{2r}$ 的二次不等式，得到
\[
 \norm{\TriH A}_{2r}\leq
 \bigl(h_r+\sqrt{h_r^2+2}\bigr)\norm A_{2r}.
\]
从 $h_2\leq1$ 出发归纳，便得到 $r=2,4,8,\ldots$ 时的统一有界性。
Schatten 类的复插值给出所有 $2\leq r<\infty$ 的有界性。
在 Hilbert--Schmidt 对偶配对下，$\TriH$ 的伴随仍为 $\TriH$，
因此对偶性又给出所有 $1<r<2$ 的有界性。

最后，
\[
 \PP=\tfrac12(I+\DD+\TriH),\qquad
 \QQ=\tfrac12(I-\DD-\TriH),
\]
故 $\PP,\QQ$ 在每个 $1<r<\infty$ 的 $S_r$ 上均一致有界。
利用有限支撑矩阵的稠密性即可延拓到整个 $S_r(H)$。
这就证明了 \eqref{eq:tp}，并完成定理 \ref{thm:bounded} 和
\ref{thm:sharp} 中所用的三角截断有界性论证。

\begin{thebibliography}{9}
\bibitem{Power}
S. Power,
\emph{Commutators with the triangular projection and Hankel forms on nest algebras},
J. London Math. Soc. (2) \textbf{32} (1985), 272--282.
\href{https://doi.org/10.1112/jlms/s2-32.2.272}{doi:10.1112/jlms/s2-32.2.272}.
本文使用其第 277 页的距离公式，并与定理 2.1、2.3
以及第 279 页的 Hankel 分量分解比较。
\end{thebibliography}

\end{document}
